\chapter{量化 C++ 求职与入职}

\begin{introduction}
  \item 把语言知识组织为代码、性能和系统设计能力
  \item 用证据描述项目，不堆砌技术名词
  \item 建立面试复盘与入职学习循环
  \item 识别不同量化岗位对 C++ 深度的差异
\end{introduction}

\chaptermeta{完成贯穿项目并能解释其测试与限制。}{Python 经历是优势，但面试回答必须补充 C++ 的生命周期、成本模型和生产边界。}{准备材料若只有形容词没有命令、测试或数字，应删除形容词并补可复现证据。}

\section{先识别岗位}

Quant Researcher 更重模型、实验和数据，但需要理解原生库与性能边界；Research Engineer 连接研究与生产，重视可复现数据管道、数值正确性和 Python/C++ 接口；Quant Developer 构建回测、执行、风险和基础设施；低延迟工程师进一步要求操作系统、网络、硬件计数器和并发内存模型。

不要用同一份准备清单覆盖所有岗位。阅读职位描述时，把要求映射为语言、算法、系统、金融领域和协作五列，再为每项准备“做过什么、证据是什么、还不知道什么”。

\section{语言面试主线}

高频问题不是孤立关键字，而是成本与契约：

\begin{itemize}
  \item 对象生命周期、RAII、复制/移动与悬空；
  \item 值语义、引用、\texttt{const} 和智能指针选择；
  \item 容器布局、迭代器失效和复杂度；
  \item 模板、concepts、虚分派与 ABI；
  \item 异常安全、错误边界和 UB；
  \item 缓存、分配、benchmark 与 profiler；
  \item 数据竞争、锁、原子、队列和停止协议。
\end{itemize}

回答采用“定义—何时使用—代价—失败例子”。例如 shared pointer：引用计数共享生命周期；用于没有自然单一所有者的对象；成本包含控制块和原子计数；环会泄漏且不提供对象内部线程安全。

\section{代码题}

先澄清输入规模、所有权、错误和复杂度，再写小而正确的实现。选择 \texttt{vector} 作为默认序列，使用标准算法但不炫技。提交前检查空输入、整数溢出、索引边界、迭代器失效、比较器严格弱序和异常路径。

性能题不要立刻写 SIMD 或无锁队列。先定义指标和工作负载，给出基线表示，指出潜在缓存/分配成本，再说明如何测量。面试官通常更看重推理闭环。

\section{系统设计题}

量化系统设计可以沿数据生命线回答：来源、时间与排序、验证、存储、消费、状态、错误、回放、观测和容量。实时行情系统必须回答背压与断线恢复；回测平台必须回答数据版本、确定性、资源隔离与结果可追溯；订单系统必须回答幂等、状态机、风险门和审计。

明确一致性需求。交易前风险可能要求同步拒绝，研究统计可以最终一致；市场事件可能按交易所序列排序，而跨市场只有定义好的合并规则。不要用“Kafka + Redis + 微服务”替代需求分析。

\section{项目如何写进简历}

弱表述：“使用 C++、CMake 和多线程开发高性能回测框架。”它没有规模、行为和证据。

更可信的结构是：动作 + 系统边界 + 验证 + 测量。例如：

\begin{quote}
用 C++20 构建事件驱动回测内核，将 CSV 验证、策略、撮合和组合划分为四个公开测试 seam；以手算账本验证端到端收益，并建立 AoS/SoA 可重复基准，明确硬件与工作负载限制。
\end{quote}

只有真实测量后才能加入吞吐或延迟数字，并保留复现实验命令。不要把教学模型称为生产级交易系统。

\section{行为面试}

准备 STAR 故事时优先选择有冲突和证据的工程事件：发现假绿色测试、反驳没有数据的优化、处理数据质量事故、在截止期内缩小范围、与研究者澄清可交易假设。说明你如何改变流程，而不仅是修了一行代码。

\section{入职后的前三个月}

前 30 天建立系统地图：构建、测试、部署、数据血缘、关键指标和术语；31--60 天完成小型端到端变更，理解审查与发布；61--90 天在证据支持下改善一个可靠性或性能问题。优先学习团队代码库的约定，不要把个人偏好当现代化改造。

\begin{interviewcheck}
用两分钟解释本书项目；回答“为什么不用 shared pointer 管理所有事件”；设计行情队列的背压；给出一个带基线和限制的性能结论；说出当前项目最重要的三个非目标。
\end{interviewcheck}

\section{练习}

\begin{enumerate}
  \item 为目标职位做五列能力矩阵，并为每项附一条证据。
  \item 录制五分钟项目讲解，删掉无法证明的形容词。
  \item 完成一次 45 分钟代码题，复盘正确性、复杂度和沟通。
  \item 为行情处理或回测平台画系统图，标注背压、恢复和观测点。
\end{enumerate}

\section{本章小结}

求职准备应把 C++ 特性转化为生命周期、成本、正确性和系统边界的推理。项目陈述依靠可复现证据，系统设计沿数据生命线展开，入职后则从理解现有约束开始。
